\chapter{Guía de despliegue y reproducibilidad}
\label{anexo:despliegue}

Este anexo recoge las instrucciones concretas para reproducir el
pipeline y desplegar la interfaz en un equipo nuevo. Se documentan los
requisitos previos, la instalación de dependencias, la configuración de
la clave API de Groq, la ejecución del pipeline y el arranque de la
aplicación web.

\section*{Requisitos previos}
\addcontentsline{toc}{section}{Requisitos previos}

\begin{itemize}
  \item Python 3.12 o superior (probado con 3.12).
  \item Sistema operativo con soporte para las dependencias científicas
  estándar: macOS, Linux o Windows (subsistema WSL2 recomendado).
  \item Aproximadamente 3\,GB de espacio libre en disco para las
  cohortes descargadas y los artefactos intermedios.
  \item Cuenta de Groq (\url{https://console.groq.com/}) con clave API
  activa; la tier gratuita alcanza para la curación de las 11 cohortes
  y para varias sesiones diarias de asistente conversacional.
  \item Conexión a Internet para la descarga de las cohortes GEO y las
  llamadas a la API de Groq.
\end{itemize}

\section*{Instalación de dependencias}
\addcontentsline{toc}{section}{Instalación de dependencias}

Desde la raíz del repositorio:

\begin{lstlisting}[style=bash, caption={Instalación en entorno virtual.}]
python -m venv .venv
source .venv/bin/activate   # o .venv\Scripts\activate en Windows
pip install -r requirements.txt
\end{lstlisting}

El fichero \texttt{requirements.txt} incluye:
\texttt{pandas}, \texttt{numpy}, \texttt{scipy}, \texttt{scikit-learn},
\texttt{statsmodels}, \texttt{GEOparse}, \texttt{groq},
\texttt{streamlit}, \texttt{python-dotenv}, \texttt{matplotlib},
\texttt{seaborn} y \texttt{pytest} (esta última sólo para tests).

\section*{Configuración de la clave API}
\addcontentsline{toc}{section}{Configuración de la clave API}

Crear un fichero \texttt{.env} en la raíz del repositorio con el
contenido:

\begin{lstlisting}[caption={Fichero .env de configuración.}]
GROQ_API_KEY=gsk_XXXXXXXXXXXXXXXXXXXXXXXXXXXXXXXXXXXXXXXXX
\end{lstlisting}

El fichero \texttt{.env} está incluido en \texttt{.gitignore} para
evitar filtrar la clave al repositorio público. Sin esta clave el
pipeline no puede curar los metadatos (paso 2), pero puede reejecutarse
si ya existe un \texttt{AUDITORIA\_COHORTES.csv} previo con las
etiquetas asignadas.

\section*{Ejecución del pipeline completo}
\addcontentsline{toc}{section}{Ejecución del pipeline completo}

Los pasos en orden lógico, ejecutables desde la raíz:

\begin{lstlisting}[style=bash, caption={Reproducción end-to-end.}]
# Fase 1: descarga
python agentes/paso1_descarga.py

# Fase 2: curacion y normalizacion
python agentes/paso2_metadata.py

# Fase 3: analisis diferencial por cohorte
python agentes/paso3_diferencial.py

# Fase 4-6: informes y figuras por cohorte
python agentes/paso4_informe.py
python agentes/paso5_graficos.py
python agentes/paso7_orquestador.py

# Fase 7: modelo transversal (LODO, subtipo, firma)
python agentes/paso13_subtipo_lodo.py
python agentes/paso14_auditoria_datos.py
python agentes/paso15_lodo_honesto.py
python agentes/paso16_composicion_vs_biologia.py
python agentes/paso17_falacia_folds.py
python agentes/paso18_subtipo_casos_dificiles.py
python agentes/paso19_firma_validada.py

# Figuras globales
python agentes/generar_figuras_auditoria.py
\end{lstlisting}

La ejecución completa dura aproximadamente 25--40 minutos en un equipo
de escritorio moderno (M1 o Intel i7), dominadas por la descarga de las
cohortes (fase 1) y por la curación LLM (fase 2). Los pasos 13-19
tardan típicamente entre 3 y 6 minutos.

\section*{Ejecución selectiva}
\addcontentsline{toc}{section}{Ejecución selectiva}

Cada paso persiste sus salidas en disco, por lo que no es necesario
reejecutar todo el pipeline cuando se cambia un solo módulo. Los
escenarios más frecuentes son:

\begin{description}
  \item[Añadir una cohorte:] editar \texttt{datasets.txt}, reejecutar
  fases 1 y 2 (\texttt{paso1}, \texttt{paso2}), y a partir de ahí toda
  la fase 7 y las figuras.
  \item[Cambiar el clasificador:] reejecutar sólo la fase 7 desde el
  paso que corresponda (\texttt{paso15\_lodo\_honesto.py} para la tarea
  tumor-vs-sano; \texttt{paso13\_subtipo\_lodo.py} para el subtipo).
  \item[Recomputar la firma:] reejecutar únicamente
  \texttt{paso19\_firma\_validada.py}.
\end{description}

\section*{Interfaz web}
\addcontentsline{toc}{section}{Interfaz web}

Una vez ejecutado el pipeline (o al menos los pasos 14, 15, 16, 17, 18,
19), la interfaz se arranca con:

\begin{lstlisting}[style=bash, caption={Arranque de la interfaz.}]
streamlit run agentes/paso12_web_chatbot.py
\end{lstlisting}

La aplicación abre por defecto en \url{http://localhost:8501}. Para
desplegarla en un servidor accesible desde red, añadir los flags:

\begin{lstlisting}[style=bash]
streamlit run agentes/paso12_web_chatbot.py \
    --server.address=0.0.0.0 \
    --server.port=8080 \
    --server.headless=true
\end{lstlisting}

\section*{Tests automáticos}
\addcontentsline{toc}{section}{Tests automáticos}

La suite mínima de tests se ejecuta con:

\begin{lstlisting}[style=bash]
pytest tests/
\end{lstlisting}

Los tests cubren el alineamiento muestra-etiqueta, que es la operación
más crítica del pipeline. Cuando no hay cohortes descargadas localmente
(por ejemplo en el runner de GitHub Actions), los tests se saltan con
\texttt{pytest.skip} en lugar de fallar.

El flujo de integración continua está configurado en
\texttt{.github/workflows/tests.yml} y se ejecuta automáticamente en
cada \emph{push} y \emph{pull request} a la rama principal.

\section*{Integración con LaTeX (esta memoria)}
\addcontentsline{toc}{section}{Integración con LaTeX (esta memoria)}

La memoria del TFM (fuente en \texttt{memoria/}) se compila con la
siguiente secuencia:

\begin{lstlisting}[style=bash, caption={Compilación de la memoria.}]
cd memoria
latexmk -pdf main.tex
\end{lstlisting}

O manualmente:

\begin{lstlisting}[style=bash]
cd memoria
pdflatex main.tex
bibtex main
pdflatex main.tex
pdflatex main.tex
\end{lstlisting}

La memoria depende de los siguientes paquetes LaTeX (incluidos en las
distribuciones estándar TeX Live y MacTeX):
\texttt{babel-spanish}, \texttt{lmodern}, \texttt{geometry},
\texttt{setspace}, \texttt{fancyhdr}, \texttt{graphicx}, \texttt{float},
\texttt{booktabs}, \texttt{longtable}, \texttt{multirow},
\texttt{subcaption}, \texttt{xcolor}, \texttt{listings},
\texttt{amsmath}, \texttt{mathtools}, \texttt{hyperref},
\texttt{natbib}, \texttt{titlesec}, \texttt{appendix}.

Todas las figuras y tablas empleadas en la memoria son los artefactos
generados por el pipeline y residen en \texttt{memoria/figuras/}. La
bibliografía está en \texttt{memoria/bibliografia.bib} en formato
BibTeX.

\section*{Resolución de problemas frecuentes}
\addcontentsline{toc}{section}{Resolución de problemas frecuentes}

\begin{description}
  \item[La descarga de una cohorte falla.] GEO limita el número de
  peticiones simultáneas. Reintentar tras unos minutos suele funcionar.
  Si persiste, comprobar la conectividad y el servicio en
  \url{https://www.ncbi.nlm.nih.gov/geo/}.
  \item[El LLM devuelve JSON inválido.] El pipeline registra estas
  muestras como \texttt{ambiguo} y sigue adelante. Si la tasa de éxito
  cae por debajo del 40\,\%, la cohorte queda excluida por criterios de
  inclusión y no rompe el análisis.
  \item[Streamlit no encuentra la clave API.] Verificar que el fichero
  \texttt{.env} está en la raíz del proyecto y que la clave empieza por
  \texttt{gsk\_}. Sin clave, el asistente queda deshabilitado pero el
  resto de la interfaz funciona.
  \item[Los pantallazos del asistente no se generan.] Es normal en el
  runner de CI si no hay clave API. Ejecutar en local con clave para
  obtener los pantallazos actualizados.
\end{description}
